\scanpage{292}
\mspage{33}
On va d'abord regarder un foncteur quelconque
\[
 f\colon \unsure{\mathbb{\Delta}|\Delta}_i^\varepsilon\longrightarrow I_{\unsure{\tau|t}'}
 \qquad(\varepsilon\in\{\pm1\deleted{\unsure{-}}\}).
\]
Supposons pour \deleted{exemple} \inserted{fixer les idées} que $\varepsilon=+1$, i.e. que \unsure{le|la}
\unsure{premier mb|premier int} soit un intervalle \deleted{\unsure{dire}} critique
\unsure{direct}\ldots{} On trouve que $f$ \deleted{applique} se factorise
par un sous-intervalle critique de $I_{\tau'}$,
\unsure{soit} $\mathbb{\Delta}_j^{\varepsilon'}$.
% slanted
\margin{\uline{Lemme} Si $I_{\tau'}$ est un intervalle, \deleted{\unsure{deux|tout}}\unsure{une|toute} partie $J$ de $I_{\tau'}$ telle que \deleted{\unsure{pour}} pour deux éléments de $J$, il existe tjrs une flèche de l'un à l'autre, est \unsure{contenue|contenu} dans un intervalle critique (unique si $\operatorname{Card}J\geq2$)}
Il est donc défini par une application
croissante de $\Delta_i^\varepsilon$ dans $\Delta_j^{\varepsilon'}$ (si $\varepsilon,\varepsilon'$ de $=$ signe),
une application décroissante dans \unsure{le} sens inverse.

Mais si on exige que $f$ soit croissant, alors
il faut que \uline{ou bien} \deleted{\unsure{$\varepsilon=\varepsilon'$, ou}} $f$ soit \unsure{constante|constant}
(auquel cas \deleted{\unsure{$I_{\tau'}$}}, si sa valeur n'est pas une extrémité
de $I_{\tau'}$, il y a exactement \uline{deux} intervalles
critiques de $I_{\tau'}$ \deleted{qui} \deleted{\unsure{se}} par lesquels se factorise $f$),
ou bien $f$ non constante et $\varepsilon=\varepsilon'$ (et en tous cas,
si $f$ non constante, alors l'intervalle critique par
lequel se factorise $f$ est unique). Pour
résumer :

\uline{Prop 7} Soient $I_\tau,I_{\tau'}$ deux intervalles orientés,
\[
 f\colon I_\tau\unsure{\mathbin{-}|\longrightarrow}I_{\tau'}
\]
un foncteur. Alors pour \deleted{tout} intervalle critique
$J\subset I_\tau$ tel que $f\mid J$ soit non constant, il existe
un unique intervalle critique $J'\subset I_{\tau'}$ tel que
$f\mid J$ se factorise par $J'$ i.e. tel que $f(J)\subset J'$.
Si $f$ est croissant, alors $J$ et $J'$ sont de $=$ signe,
% archive sheet number: 291
